\chapter{Statistiques et lecture critique}
\label{corriges:10}
\section*{Exercice 10.1 -- Tous les indicateurs}
$N=16$. Les notes au moins égales à $14$ regroupent $4+1=5$ élèves, soit $5/16=31{,}25\,\%$.
La somme vaut $8\times2+10\times3+12\times6+14\times4+16=190$, donc $\bar x=190/16=11{,}875$.
Les rangs $6$ à $11$ valent $12$, donc les rangs centraux $8,9$ valent $12$ : médiane $12$.
Étendue : $16-8=8$.

\section*{Exercice 10.2 -- Choisir et construire un graphique}
Les trois classes ont la même amplitude $10$. Dessiner trois rectangles adjacents de largeurs $10$ et de hauteurs $5,9,6$. Le total vaut $20$. L'angle de la deuxième classe vaut $360\times9/20=162\degre$.
La moyenne exacte est inconnue : les positions dans les classes ne sont pas données. Avec les centres $5,15,25$, on peut estimer $(5\times5+9\times15+6\times25)/20=15{,}5$ minutes, sans affirmer l'exactitude.

\section*{Exercice 10.3 -- Moyenner des moyennes}
Le total des notes vaut $12\times10+18\times15=390$ pour $30$ élèves. La moyenne commune est $390/30=13$. La réponse $(10+15)/2=12{,}5$ donne le même poids aux deux groupes alors qu'ils n'ont pas le même effectif.

\section*{Exercice 10.4 -- Valeur extrême et lecture critique}
Première série : moyenne $4$, médiane $4$, étendue $4$.
Deuxième : moyenne $40/5=8$, médiane $4$, étendue $24$.
La médiane est inchangée. Une origine à $20$ peut faire paraître l'écart entre $21$ et $22$ très grand visuellement, surtout avec des barres. Il faut lire l'échelle : la hausse réelle est $1/21\approx4{,}76\,\%$, pas un doublement.
